\documentclass[11pt,letterpaper]{article}
\usepackage[T1]{fontenc}
\usepackage[utf8]{inputenc}
\usepackage{lmodern}
\usepackage[margin=0.9in,headheight=15pt,headsep=18pt]{geometry}
\usepackage{amsmath,amssymb,mathtools,mathrsfs}
\usepackage{microtype}
\usepackage{booktabs,longtable,array}
\usepackage{graphicx}
\usepackage[section]{placeins}
\usepackage{xcolor}
\usepackage{enumitem}
\usepackage{fancyhdr}
\usepackage{titlesec}
\usepackage{tocloft}
\setlength{\cftbeforesecskip}{3pt}
\renewcommand{\cftsecfont}{\small}
\renewcommand{\cftsecpagefont}{\small}
\usepackage{needspace}
\usepackage{xurl}
\usepackage{hyperref}
\hypersetup{hidelinks,pdftitle={Komorebi Large Signal Model: Foundations 0.8.0},pdfauthor={Komorebi P.}}
\setlength{\parindent}{0pt}
\setlength{\parskip}{0.52em}
\setlength{\emergencystretch}{2em}
\setcounter{secnumdepth}{-1}
\setcounter{tocdepth}{1}
\titleformat{\section}{\Large\bfseries}{\thesection}{0em}{}
\titlespacing*{\section}{0pt}{1.7em}{0.7em}
\titleformat{\subsection}{\large\bfseries}{\thesubsection}{0em}{}
\titlespacing*{\subsection}{0pt}{1.15em}{0.5em}
\clubpenalty=10000
\widowpenalty=10000
\displaywidowpenalty=10000
\allowdisplaybreaks[1]
\providecommand{\tightlist}{\setlength{\itemsep}{0pt}\setlength{\parskip}{0pt}}
\providecommand{\passthrough}[1]{#1}
\pagestyle{fancy}
\fancyhf{}
\fancyhead[L]{\small\textsc{Komorebi Large Signal Model}}
\fancyhead[R]{\small Foundations 0.8.0}
\fancyfoot[L]{\small Komorebi P.}
\fancyfoot[R]{\small\thepage}
\renewcommand{\headrulewidth}{0.3pt}
\begin{document}
\hypersetup{pageanchor=false}
\begin{titlepage}
\vspace*{0.5in}
{\small\textsc{Operation Komorebi}\par}
\vspace{0.72in}
{\fontsize{33}{39}\selectfont\bfseries Large Signal Model\par}
\vspace{0.3in}
{\fontsize{21}{27}\selectfont Data-supported domains\\and structured uncertainty\par}
\vspace{0.45in}
\rule{\textwidth}{0.45pt}\par
\vspace{0.3in}
\begin{minipage}{0.95\textwidth}
\large
Eight statements with proofs.\\Confidence from dependent observations.\\Smaller repairs and explicit prediction bounds.
\end{minipage}
\vfill
{\Large Komorebi P.\par}
\vspace{0.2in}
{\normalsize Foundations edition 0.8.0\par}
{\normalsize Companion to PR-LASM software release 2.3.0\par}
{\normalsize 9 October 2026\par}
\end{titlepage}
\hypersetup{pageanchor=true}
\begingroup
\footnotesize
\setlength{\cftbeforesecskip}{1pt}
\tableofcontents
\endgroup
\newpage

\section{Part VIII. Data-supported domains and smaller
repairs}\label{part-viii.-data-supported-domains-and-smaller-repairs}

The preceding block certified a perturbed operator family after mixing
it with a center reset. Its all-horizon validity guarantee was exact,
but its error bound became trivial at horizon eight. It also kept
deterministic coefficient boxes separate from statistical confidence
sets. This block addresses both limitations in declared sectors rather
than weakening either distinction.

First, the same fitted family and previously proposed geometry admit a
smaller, independently checked repair. Second, a self-normalized
martingale argument produces operator uncertainty sets from a dependent
trajectory when the predictive state is an accessible, fixed observable.
A nonlinear three-coordinate example connects those sets to a curved
domain without a reset correction. Its geometry and feature dictionary
are specified in advance. This is not a solution to discovering an
unknown hidden state from the same data.

The sequential-regression result is a classical self-normalized
concentration theorem {[}R23{]}, proved here in the form used by the
signal model. The contraction argument uses standard Markov-kernel and
total-variation ideas {[}R24{]}. The new block's contribution is the
explicit sequence of assumptions, finite witnesses, structured
uncertainty, observation-measure distinctions and executable examples
connecting these tools to the realization calculus. Proposition counts
are not historical-novelty counts.

\newpage

\section{91. A smaller repair can be bracketed rather than
guessed}\label{a-smaller-repair-can-be-bracketed-rather-than-guessed}

Keep the fitted operators, proposed transform and initial state from
Section 88 unchanged. In Lorentz coordinates use the center-reset
operator \(N=e_0\lambda\), of constant density one under \(du/2\). Write

\[
F_\eta=(1-\eta)\widehat F+\eta N,
\qquad 0\leq\eta\leq1.
\]

Fix a required margin \(0\leq\kappa\leq1\). Feasibility means that every
action, observation and normalized input has margin at least \(\kappa\).
Integrated mass and initial membership remain separate requirements.

\Needspace{5\baselineskip}
\textbf{Proposition 80. Convex repair bracketing.} The feasible repair
weights form a closed interval \([\eta_*,1]\). If a cone-positive
functional \(\ell=(1,-v^{\mathsf T})\), with \(\|v\|_2\leq1\), and a
normalized input \(z\) give \(c=\ell\widehat F_a(u)z<\kappa\), then
every feasible weight satisfies

\[
\boxed{\eta\geq\frac{\kappa-c}{1-c}.}
\]

A complete certificate at a weight \(\eta_U\) gives an upper bound on
\(\eta_*\). A negative Bernstein coefficient without a negative
evaluated point does not give a lower bound.

\textbf{Proof.} For a fixed input and observation, output mass minus
spatial norm is a concave continuous function of the operator, hence of
\(\eta\). Its superlevel set is closed and convex. Intersecting these
sets preserves those properties. The reset has margin one, so the
intersection contains one and is an interval ending there. If \(F_\eta\)
has margin at least \(\kappa\), every such functional satisfies
\(\ell F_\eta z\geq\kappa\). Since \(\ell Nz=1\), this requires
\((1-\eta)c+\eta\geq\kappa\). Here \(1-c>0\), and solving gives the
bound. A complete positive witness proves feasibility. A negative
coefficient only says that a particular sufficient polynomial bound
failed. \(\square\)

\subsection{91.1 The unchanged predecessor
example}\label{the-unchanged-predecessor-example}

Use the previous transform with center \((1/2,1/2)\) and triangular
factor

\[
L=\begin{pmatrix}123/37&-2\\0&109/41\end{pmatrix}.
\]

At the direct action, observation \(u=-1\), normalized input \((1,1,0)\)
and functional \(\ell=(1,7/25,-24/25)\), exact evaluation gives

\[
c=-\frac{17590411}{2332387500}.
\]

For the required margin \(\kappa=1/64\), this proves

\[
\frac{864543451}{37599646576}
\leq\eta_*\leq\frac1{32}.
\]

The lower endpoint is approximately \(0.02299\). The upper endpoint is
not a sampled estimate. A complete exact Bernstein certificate proves
that \(\eta=1/32\) works. The weight \(1/64\) fails the requested
\(1/64\) margin at the displayed separating functional. This does not by
itself prove that its nominal model is probability-invalid.

The new certificate uses 194 leaves and 3,058 recomputed coefficients. A
separate comparison-domain certificate uses 185 leaves and 2,833
coefficients. Both retain the same degree-two coefficient box of radius
\(1/4096\) in their respective declared Lorentz coordinates, subject to
exact integrated mass. The more general mass and spatial output budgets
remain \(1/128\) each. Their sum equals the certified margin, so that
larger output-error family has a nonnegative lower margin, not a claimed
strictly positive residual margin.

The reset weight has fallen from \(1/8\) to \(1/32\), while the
deterministic uncertainty box is unchanged. No global optimization over
geometry or repair operators is claimed. The current bracket concerns
this fixed geometry, this reset path and this required margin.

The comparison domain is checked again before invoking the law theorem.
Keeping the same raw discrepancy budget \(241/40000\) from Section 89
gives

\[
\varepsilon_{\mathrm{repair}}
\leq\frac{31}{32}\frac{241}{40000}+\frac1{16}
=\frac{87471}{1280000}.
\]

Consequently, under a common policy and identical initial states,

\[
\boxed{
\|P_T-\widetilde P_T\|_{\mathrm{TV}}
\leq\min\left(1,\frac{87471T}{2560000}\right).
}
\]

The first integer horizon where this upper bound reaches one is now 30
rather than eight. The bound remains conservative. The repaired
uncertainty tube is still not a confidence set for the unmodified
generating process. This improvement is deterministic certificate
refinement, not a new statistical experiment.

\newpage

\section{92. Confidence from dependent
observations}\label{confidence-from-dependent-observations}

Let \(\mathcal F_{t-1}\) contain all information available before
observation \(t\), including its already selected action. Let
\(\phi_t\in\mathbb R^d\) be measurable with respect to that information.
Assume

\[
Y_t=\theta_*^{\mathsf T}\phi_t+\xi_t,
\qquad
\mathbb E[e^{h\xi_t}\mid\mathcal F_{t-1}]
\leq e^{h^2\sigma^2/2}
\]

for every real \(h\), with \(\sigma>0\), and a declared bound
\(\|\theta_*\|_2\leq S\). With ridge \(\gamma>0\), set

\[
V_t=\gamma I+\sum_{s\leq t}\phi_s\phi_s^{\mathsf T},
\qquad
\widehat\theta_t=V_t^{-1}\sum_{s\leq t}\phi_sY_s.
\]

\Needspace{5\baselineskip}
\textbf{Proposition 81. Time-uniform regression confidence.} For
\(0<\delta<1\), with probability at least \(1-\delta\), simultaneously
for every finite \(t\),

\[
\boxed{
\|\widehat\theta_t-\theta_*\|_{V_t}
\leq
\sigma\sqrt{
\log\frac{\det V_t}{\det(\gamma I)}+2\log\frac1\delta
}
+\sqrt\gamma S.
}
\]

No independence of successive observations or of the regressors is
required. The conditional mean, conditional noise and predictability
hypotheses are required. The confidence set remains valid at a stopping
time selected from the data.

\textbf{Proof.} Put \(S_t=\sum_{s\leq t}\phi_s\xi_s\). For every fixed
\(q\), the conditional moment bound makes

\[
\exp\left(q^{\mathsf T}S_t/\sigma
-\tfrac12q^{\mathsf T}(V_t-\gamma I)q\right)
\]

a nonnegative supermartingale. Integrate against the fixed Gaussian law
\(q\sim N(0,(\gamma I)^{-1})\). Conditional Tonelli and completing the
square give the nonnegative supermartingale

\[
M_t=
\left(\frac{\det(\gamma I)}{\det V_t}\right)^{1/2}
\exp\left(\frac{\|S_t\|_{V_t^{-1}}^2}{2\sigma^2}\right),
\qquad M_0=1.
\]

For completeness, stop at the first crossing of \(1/\delta\), then at
any deterministic finite horizon. The supermartingale expectation is at
most one, so the crossing probability up to that horizon is at most
\(\delta\). Increasing the horizon proves the same bound for any finite
crossing. Off this event, taking logarithms bounds
\(\|S_t\|_{V_t^{-1}}\) by the square-root term. Finally,

\[
\widehat\theta_t-\theta_*
=V_t^{-1}(S_t-\gamma\theta_*).
\]

The triangle inequality and \(V_t\succeq\gamma I\) bound the ridge term
by \(\sqrt\gamma S\). The simultaneous event already covers every
selected finite stopping time. \(\square\)

This is the self-normalized mixture argument in {[}R23{]}, not a new
concentration principle. A bounded response in an interval of length
\(b-a\) has conditionally sub-Gaussian centered noise with
\(\sigma=(b-a)/2\). One direct proof differentiates its tilted log
moment-generating function twice. The second derivative is a variance
bounded by \((b-a)^2/4\). Integrating from zero gives the stated moment
bound.

The determinant term accounts for the observed design, but does not
create excitation. If \(\phi_t\) never probes a direction, the set need
not shrink in that direction. Regressors extracted after inspecting the
same response are not predictable. An estimated hidden state cannot be
substituted for \(\phi_t\) while silently retaining the theorem's true
conditional-mean equation.

\section{93. From observable moment regression to operator
sets}\label{from-observable-moment-regression-to-operator-sets}

There is a direct connection to the unnormalized instrument. Assume that
the normalized predictive state \(Z_t\) is an exact, fixed, measurable
function of the available observation history. For action \(a\), suppose

\[
P(U_t\in du\mid\mathcal F_{t-1})
=\lambda B_a(u)Z_t\,\mu(du),
\qquad
Z_{t+1}=\frac{B_a(U_t)Z_t}{\lambda B_a(U_t)Z_t}
\]

on realized positive-density branches. Crucially, \(Z_{t+1}\) must also
be available from the observations using the specified state map. This
is an accessibility assumption, not a consequence of successful
regression.

Let \(B_a(u)=\sum_{k=0}^{D}\psi_k(u)C_{a,k}\). Choose functions \(g_j\)
so that the known matrix \(W_{jk}=\int g_j(u)\psi_k(u)\,\mu(du)\) is
invertible.

\Needspace{5\baselineskip}
\textbf{Proposition 82. Sequential observable-moment operator sets.}

Define \(M_{a,j}=\int g_j(u)B_a(u)\,\mu(du)\). Then

\[
\boxed{
\mathbb E[g_j(U_t)Z_{t+1}\mid\mathcal F_{t-1}]
=M_{a,j}Z_t,
\qquad
C_{a,k}=\sum_j(W^{-1})_{kj}M_{a,j}.
}
\]

Under the rowwise noise and parameter-bound assumptions of Proposition
81, simultaneous confidence ellipsoids for the rows of \(M_{a,j}\) give
a confidence set for every \(C_{a,k}\) by this known linear map. If row
\(i\) of \(\widehat M_{a,j}-M_{a,j}\) has \(V_{a,t}\)-norm at most
\(\beta_{a,j,i}\), then for any direction \(x\),

\[
\left|e_i^{\mathsf T}(\widehat C_{a,k}-C_{a,k})x\right|
\leq
\sum_j |(W^{-1})_{kj}|\,
\beta_{a,j,i}\,\|x\|_{V_{a,t}^{-1}}.
\]

\textbf{Proof.} In the conditional integral, multiply the normalized
state by the observation density. The denominator cancels, giving
\(\int g_j(u)B_a(u)Z_t\,\mu(du)\). The polynomial expansion gives
\(M_{a,j}=\sum_kW_{jk}C_{a,k}\) and hence the inverse formula. Apply
Proposition 81 to each row, using the predictable regressor \(Z_t\) on
visits to action \(a\), or its action-gated version. A union bound over
the predeclared finite rows and moments gives simultaneous coverage.
Independence between these response coordinates is unnecessary. Weighted
Cauchy--Schwarz and the triangle inequality give the displayed
directional bound. \(\square\)

Exact mass constraints can be imposed by intersecting the confidence set
with the known mass affine subspace. This preserves coverage because the
true model is assumed to obey that law. Arbitrarily shrinking a set
around a convenient corrected operator does not have the same
justification. Data-selected geometry can be checked against the whole
resulting set using Proposition 78.

The theorem is conditional on a correct finite observation basis,
accessible predictive state and parameter bounds. It does not obtain a
hidden predictive state from a single dependent path without additional
assumptions. Unknown dictionary discovery, state-estimation error and
polynomial misspecification remain distinct problems. The following
example makes the state accessible by construction rather than
concealing that requirement.

\newpage

\section{94. A nonlinear observable state with a curved moment
domain}\label{a-nonlinear-observable-state-with-a-curved-moment-domain}

For \(u\in[-1,1]\), put

\[
v(u)=(1,u,u^2)^{\mathsf T},
\qquad L_2(u)=\frac{3u^2-1}{2}.
\]

The normalized states are moments of a scalar signal:

\[
K=\{(1,x,y):x^2\leq y\leq1\}.
\]

Their cone has the finite curved description

\[
C=\left\{(m,x,y):
\begin{pmatrix}m&x\\x&y\end{pmatrix}\succeq0,
\quad m-y\geq0\right\},
\qquad\lambda=(1,0,0).
\]

Assume that the declared probability measure \(\mu\) on \([-1,1]\) has
moments zero at orders one and three, \(1/3\) at order two and \(1/5\)
at order four. Define a two-parameter degree-four operator density

\[
\boxed{
B_{a,b}(u)=v(u)\,(1,au,bL_2(u)).
}
\]

Here \(a,b\) are parameters, not action labels. This example is
autonomous. Its dependence on the observed signal is nonlinear because
the state includes its square.

\Needspace{5\baselineskip}
\textbf{Proposition 83. Structured curved-domain certificate.} The set
\(K\) is the convex hull of all \(v(u)\). If \(|a|+|b|\leq1-\rho\) for
\(0\leq\rho\leq1\), then \(B_{a,b}(u)\) preserves \(C\) and has exact
integrated mass. On a normalized state, its density is

\[
p_{a,b}(u\mid x,y)=1+a x u+b yL_2(u)\geq\rho.
\]

Every positive-density posterior is exactly \(v(u)\). In particular,
starting from \(v(s)\) produces the observable Markov law

\[
p_{a,b}(u\mid s)=1+a s u+b s^2L_2(u).
\]

Its next conditional moments satisfy

\[
\boxed{
\mathbb E[U_t\mid U_{t-1}=s]=\frac a3s,
\qquad
\mathbb E[L_2(U_t)\mid U_{t-1}=s]=\frac b5s^2.
}
\]

\textbf{Proof.} Every distribution on \([-1,1]\) has \(x^2\leq y\leq1\).
Conversely, fix \(|x|<1\). A point mass at \(x\) gives \((x,x^2)\),
while the endpoint mixture with mean \(x\) gives \((x,1)\). Their convex
combinations give every intermediate \(y\). The cases \(x=\pm1\) are
immediate. Homogenizing yields the matrix and scalar inequalities, which
also force the only mass-zero cone element to be zero.

On \(K\), \(|x|\leq1\), \(0\leq y\leq1\), \(|u|\leq1\) and
\(|L_2(u)|\leq1\). Thus the density is at least \(1-|a|-|b|\). Each
operator output is that nonnegative scalar times \(v(u)\), so lies in
\(C\). The first and second reference moments imply
\(\lambda\int B_{a,b}(u)\,\mu(du)=\lambda\). Normalization cancels the
scalar density, yielding \(v(u)\). Symmetry through order three gives
\(\int uL_2(u)d\mu=0\). The second and fourth moments give
\(\int L_2(u)^2d\mu=1/5\). Integrating \(u\) and \(L_2\) against the
density proves the last two formulas. \(\square\)

The posterior ray lies on a curved boundary. Arbitrary perturbations of
its output coordinates can destroy validity, as in Proposition 72.
Parameter uncertainty inside the declared family behaves differently: it
changes the likelihood but not the posterior ray. A certified rectangle
in \((a,b)\) can therefore be robust even though a free matrix-entry
ball has no positive validity radius.

This is an example of uncertainty that preserves the model's algebraic
structure. It is not a claim that every real system obeys the
two-parameter family. The feature dictionary \(1,u,u^2\), its geometry,
the reference measure and the parameter model are declared inputs to
learning.

\section{95. A fully rational dependent-data reference
experiment}\label{a-fully-rational-dependent-data-reference-experiment}

Both the continuous uniform measure \(du/2\) and the following five-node
measure satisfy the moment assumptions:

\[
\begin{array}{c|rrrrr}
u&-1&-1/2&0&1/2&1\\\hline
\mu\{u\}&7/90&16/45&2/15&16/45&7/90.
\end{array}
\]

The reference experiment uses the \textbf{five-node measure}, not
rounded observations from the continuous model. At each step, the
conditional probability of a node is its reference weight times
\(p_{a,b}(u\mid s)\). Exact rational probabilities are sampled by
integer draws. The learner receives only the node sequence, the measure
and the declared model. The true parameters are present only in the
separate generator and verification code.

Matching four moments does not identify the two observation laws. The
continuous model gives the set of five nodes probability zero, while the
discrete model gives it probability one. Their single-observation total
variation is therefore one. The continuous theorems and the exact
discrete reference calculation are deliberately distinguished.

For a path of \(n\) transitions, use two scalar regressions,

\[
\begin{aligned}
\phi_{a,t}&=U_{t-1}/3,&Y_{a,t}&=U_t,&\sigma_a&=1,\\
\phi_{b,t}&=U_{t-1}^2/5,&Y_{b,t}&=L_2(U_t),&\sigma_b&=3/4.
\end{aligned}
\]

The noise scales follow from response ranges \([-1,1]\) and
\([-1/2,1]\). The prior parameter bounds are \(|a|,|b|\leq1\). Each
regression uses ridge one and failure probability \(1/40\), giving a
total failure budget \(1/20\) for both parameters at every time.

\Needspace{5\baselineskip}
\textbf{Proposition 84. Exact rational outer confidence intervals.}
Write \(V=1+\sum\phi_t^2\) and \(\widehat\theta=\sum\phi_tY_t/V\). Let
\(k=\lceil\log_2 V\rceil\) and \(h=\lceil\log_2(1/\delta)\rceil\). A
rational upper enclosure of

\[
\sigma\sqrt{\frac{(7/10)(k+2h)}{V}}+\sqrt{\frac1V}
\]

is an upper bound on Proposition 81's scalar radius. Rounding the lower
interval endpoint downward and the upper endpoint upward, then
intersecting with \([-1,1]\), preserves its conditional coverage. Both
intervals cover their true parameters simultaneously at every time with
probability at least \(19/20\).

\textbf{Proof.} The positive exponential series gives
\(e^{7/10}>1+7/10+(7/10)^2/2+(7/10)^3/6+(7/10)^4/24>2\). Hence
\(\log 2<7/10\), so \(\log V+2\log(1/\delta)\leq(7/10)(k+2h)\). Divide
the scalar regression bound by \(\sqrt V\). For a nonnegative rational
\(x\), the smallest integer \(j\) with \((j/10^{12})^2\geq x\) supplies
a rational upper square-root enclosure. Integer square root and an exact
final comparison compute it. Outward rounding only enlarges the
interval. Intersecting with a true prior constraint preserves coverage.
Apply the two-regression union bound on their time-uniform events.
\(\square\)

All counts, sums, radius enclosures and interval endpoints in this
experiment are rational. No floating-point logarithm is trusted in the
confidence arithmetic. An exact replay verifies the arithmetic and the
declared protocol. It does not authenticate the observations, prove that
a physical process satisfies the model, or measure the procedure's
coverage from one pseudorandom run.

\subsection{95.1 Executed result}\label{executed-result}

The seeded reference path uses \(a=2/5\), \(b=1/5\) and initial
observation zero. It was first run to 65,536 transitions and then
extended along the same seeded path to 262,144. Time-uniform coverage
permits that extension without spending another time-selection budget.
This is not a new independent coverage experiment.

At the final checkpoint, outward-rounded intervals are

\[
\boxed{
\frac{89}{250}\leq a\leq\frac{231}{500},
\qquad
\frac{29}{250}\leq b\leq\frac{29}{100}.
}
\]

Every member of this rectangle has density at least

\[
\rho=1-\frac{231}{500}-\frac{29}{100}=\frac{31}{125}.
\]

The rounded fitted center is \((409/1000,203/1000)\). It lies inside the
rectangle. No reset or other coefficient correction is applied.

At earlier recorded checkpoints the sufficient density-floor test
withheld a whole-box certificate. Its first successful recorded
checkpoint was 32,768 transitions, with floor \(1/100\). This is not a
claim that 32,768 is the smallest possible sample size, or even the
first passing individual step. Only the declared checkpoints were
evaluated. A withheld whole-box certificate does not declare the nominal
fit or true process invalid.

\newpage

\section{96. Structured uncertainty has an observable error
bound}\label{structured-uncertainty-has-an-observable-error-bound}

Consider two parameters in a certified family. For normalized
\(z=(1,x,y)\), their density difference is

\[
\Delta p(u)=\Delta a\,x u+\Delta b\,yL_2(u).
\]

The state-update vector is the same \(v(u)\) in both models after
observing \(u\). This does not make the models equivalent. Their
likelihoods and hence their predicted observation laws can differ.

\Needspace{5\baselineskip}
\textbf{Proposition 85. Parameter-to-law error without a
state-estimation term.} Let

\[
c_1=\int|u|\,d\mu,
\qquad c_2=\int|L_2(u)|\,d\mu.
\]

If \(|\Delta a|\leq d_a\) and \(|\Delta b|\leq d_b\), the conditional
observation laws on any common normalized state obey

\[
\boxed{
\|P(\cdot\mid z)-\widehat P(\cdot\mid z)\|_{\mathrm{TV}}
\leq\Delta:=\frac{c_1d_a+c_2d_b}{2}.
}
\]

The same bound holds on every realized common observation history. With
the same initial observation, complete future records of length \(T\)
have total variation at most \(\min(1,T\Delta)\). For the five-node
measure,

\[
c_1=\frac{23}{45},\qquad c_2=\frac{14}{45}.
\]

For \(du/2\), the respective constants are \(1/2\) and \(2/(3\sqrt3)\).

\textbf{Proof.} The first claim follows by integrating \(|\Delta p|/2\),
using \(|x|\leq1\) and \(y\leq1\). After the same observed value both
models have exactly \(v(u)\), so the conditional bound applies again
without accumulated coordinate-state error. Telescoping conditional
kernels gives the record bound, equivalently the common-history coupling
argument of Proposition 69. Direct summation gives the discrete
constants. For the continuous constant \(c_2\), split the integral at
\(u=\pm1/\sqrt3\) and integrate the quadratic polynomial with its
respective signs. \(\square\)

The use of \(\Delta\) here is an upper bound proved from a confidence
set and a model class. It is not the lower discrepancy bound from a
rejected residual test in Part VI. A test that does not reject cannot
supply the displayed upper bound by itself.

For the final reference rectangle and fitted center,

\[
d_a=\frac{53}{1000},\qquad d_b=\frac{87}{1000},
\qquad
\boxed{\Delta=\frac{2437}{90000}.}
\]

These bounds compare the fitted model with every member of the
rectangle. They apply to the true process on the conditional coverage
event. Even when that event fails, the rectangle's internal probability
certificate remains an exact statement about its members.

\section{97. Long-horizon marginal control is not full-record
control}\label{long-horizon-marginal-control-is-not-full-record-control}

The previous record bound grows with \(T\). A different question can
have a uniform answer: how far apart are the distributions of the signal
at a single future time? A lower density bound permits this distinction
to be made exactly.

Let \(P,Q\) be Markov kernels on a common state space. Suppose
\(P(s,\cdot)\geq\rho\mu(\cdot)\) for every \(s\), where \(\mu\) is one
probability measure and \(0<\rho\leq1\). Suppose also
\(\sup_s\|P(s,\cdot)-Q(s,\cdot)\|_{\mathrm{TV}}\leq\Delta\).

\Needspace{5\baselineskip}
\textbf{Proposition 86. Uniform marginal and rolling-window stability.}
For initial laws at distance \(e_0\), the one-time marginal laws satisfy

\[
\boxed{
e_t\leq\min\left(1,
(1-\rho)^te_0+
\frac{\Delta[1-(1-\rho)^t]}{\rho}\right).
}
\]

For \(k\) consecutive observations starting at time \(t\), their joint
laws are at distance at most \(\min(1,e_t+(k-1)\Delta)\). A uniform
small marginal bound does not in general bound the total variation of
arbitrarily long complete records away from one.

\textbf{Proof.} For \(\rho<1\), write
\(P(s,\cdot)=\rho\mu+(1-\rho)R(s,\cdot)\), where \(R\) is a Markov
kernel. Applying a Markov kernel cannot increase total variation, by
integration of bounded test functions. Thus
\(\|\nu P-\widetilde\nu P\|_{\mathrm{TV}}\leq(1-\rho)\|\nu-\widetilde\nu\|_{\mathrm{TV}}\).
Adding and subtracting \(\widetilde\nu P\) yields
\(e_{t+1}\leq(1-\rho)e_t+\Delta\). Iterate this inequality. The case
\(\rho=1\) follows directly. Applying the finite-record telescoping
argument from the initial time \(t\) gives the rolling-window statement.

For the final assertion, take two distinct independent Bernoulli laws
with parameters strictly between zero and one. Each is a Markov kernel
with a positive common minorization. At any time their marginal distance
is the fixed parameter difference. For a full record, test whether its
empirical mean lies above the midpoint between the parameters. Under
each law, Chebyshev's inequality makes the error probability tend to
zero as the record length grows. The total variation of the record laws
therefore tends to one. \(\square\)

This is the standard minorization/contraction mechanism {[}R24{]}, with
its error accounting made explicit. It is not a stability theorem for
arbitrary feedback policies that remember the full history. The
reference model is autonomous and its current scalar observation is the
Markov state. It also does not bound long products of unnormalized
likelihood ratios.

With \(e_0=0\), \(\rho=31/125\) and \(\Delta=2437/90000\), the reference
confidence rectangle gives

\[
\boxed{
\sup_{t\geq0}e_t\leq\frac{2437}{22320}
\approx0.10919.
}
\]

This is a nontrivial all-time bound on \textbf{one-time marginal
observation distributions}. The complete-record bound is still
\(\min(1,2437T/90000)\), which becomes trivial at integer horizon 37.
The model does not exchange those two different claims. In live
filtering after a shared observed value, its state is exactly known and
the next-observation law has the separate uniform upper bound
\(2437/90000\).

\newpage

\section{98. A conditional data-to-domain
theorem}\label{a-conditional-data-to-domain-theorem}

\Needspace{5\baselineskip}
\textbf{Proposition 87. Observable structured learning with a covered
domain.} Suppose the scalar process obeys the conditional law of
Proposition 83 with the declared reference measure and fixed parameters
satisfying \(|a|,|b|\leq1\). Use the predictable regressions, fixed
dictionary and confidence protocol of Proposition 84. At a data-selected
finite time, form the full outward-rounded rectangle. If its density
floor \(\rho\) is nonnegative, that rectangle has an exact all-horizon
probability certificate on \(C\). If \(\rho>0\), every chosen center
inside it has the marginal and record bounds of Propositions 85--86
relative to every member.

With probability at least \(19/20\) over the specified stochastic
observation process, these implications cover its true parameter at
every checked time, including a stopping time. Internal validity of a
verified rectangle is deterministic. Statistical coverage remains
conditional on the declared process, moment, parameter and observation
assumptions.

\textbf{Proof.} Proposition 84 provides one simultaneous event on which
both true parameters belong to all returned intervals. Proposition 83
supplies the domain certificate for the entire rectangle whenever its
sufficient density floor is nonnegative. The selected center is one of
its members, and the maximum distances to the endpoints bound each
parameter difference. Proposition 85 gives conditional and record error.
If the floor is positive, all true conditional rows have the same lower
reference component, so Proposition 86 applies to the scalar Markov
process. Selection of the stopping time does not alter the already
simultaneous event. None of these implications verifies the assumptions
from a single sample path. \(\square\)

No claim of empirical coverage follows from the saved run. It is one
seeded finite realization of the exactly declared discrete analogue. The
general confidence result is a theorem, with assumptions recorded in the
certificate. Its independent checker rederives the sufficient
statistics, rational radii, outer endpoints, density floor and error
budgets from the saved path. It rejects a changed reference measure, a
narrower claimed interval, an altered path checksum or an unsupported
claim that the assumptions were independently verified.

This is the point where probability-domain certification and data
uncertainty become connected in an executed dependent-signal example.
The family is more restrictive than an arbitrary learned operator. That
restriction is visible and does useful work: it preserves the nonlinear
posterior curve while uncertainty changes its likelihoods.

\section{99. What the block does and does not
close}\label{what-the-block-does-and-does-not-close}

Two separate improvements have been established.

The perturbed 0.7.0 operator family now admits a four-times-smaller
reset weight at the same deterministic coefficient-box radius. The
witness is checked on both the discovered and comparison domains, and
its prediction bound remains nontrivial for a longer finite horizon. Its
geometry was reused. No new empirical coverage is attached to that
example.

The new scalar example learns two unknown coefficients from dependent
observations of a fixed nonlinear state, constructs a time-uniform
confidence rectangle, and certifies all its members on a curved domain
without a reset. It also distinguishes complete-record, current-state
and rolling-window prediction bounds. This example uses a separate
declared model and does not establish a confidence set for the old 0.7.0
operator family.

The outstanding general issue is not merely positive probabilities. It
is obtaining an adequate observable representation, a correct
observation basis and sufficiently informative data in a model class
that remains computationally tractable. Theorems for a known dictionary
cannot be used as proofs that an estimated hidden dictionary is correct.
Unmodeled measurement error, nonstationary coefficients,
response-dependent feature selection, unsupported reference measures and
insufficient action coverage remain explicit failure modes.

\section{100. Release direction after the eighth
block}\label{release-direction-after-the-eighth-block}

The cumulative theory contains 87 propositions with written proofs. The
next step toward a foundations 1.0 should be a dependency and
assumptions audit of this declared scope rather than treating every
possible extension as a prerequisite. A defensible release can freeze
exact realization, transfer and closure, finite observable learning,
controlled instruments, probability-domain certificates, quantified
approximation, and the restricted sequential learning sector now worked
through.

General nonlinear dictionary discovery and hidden-state identification
from unrestricted noisy streams remain research directions. They need
not be represented as already solved, nor silently made requirements for
the finite theorems that are proved. The reference programs remain a
small mathematics laboratory. Production PR-LASM stays at 2.3.0. No new
Lean verification, public publication or license change is claimed.

\begingroup
\small
\raggedright

\section{References and dependencies}\label{references-and-dependencies}

\textbf{{[}R23{]}} Yasin Abbasi-Yadkori, David Pál and Csaba Szepesvári.
\emph{Online Least Squares Estimation with Self-Normalized Processes: An
Application to Bandit Problems}. arXiv:1102.2670, version 1, 14 February
2011. The primary manuscript's Section 2, Theorem 3, and Section 5
supply the classical martingale-mixture and adaptive-regression context.
Proposition 81 includes its own proof, retains conditional noise and
design assumptions, and does not claim a new concentration inequality.
Source checked: \path{https://arxiv.org/pdf/1102.2670}.

\textbf{{[}R24{]}} David A. Levin and Yuval Peres, with contributions by
Elizabeth L. Wilmer. \emph{Markov Chains and Mixing Times}, second
edition, 2017. Chapters 4 and 5 develop total variation and coupling.
The present minorization argument is proved directly and does not import
a finite-state assumption into its general kernel statement. The worked
reference chain is finite. Source checked: the author-hosted
second-edition text,
\path{https://pages.uoregon.edu/dlevin/MARKOV/mcmt2e.pdf}.

\textbf{Internal dependencies.} Part VII, especially Propositions 71,
73, 75, 76 and 79, supplies the robust checker and predecessor repair
comparison. Propositions 59 and 69 supply the continuous and
conditional-kernel full-law comparisons. Their source and proof passages
are preserved unchanged. The independent source archive is
\path{Komorebi_Large_Signal_Model_Foundations_0.7.0.zip}.

\endgroup
\end{document}
